\section*{Abstract}

Converting relational data into an RDF knowledge graph (KG) is straightforward with a direct mapping, but the result mirrors how the data is stored rather than what it describes. Turning it into a semantically modeled KG, with appropriate entities, meaningful classes and predicates, and explicit relationships, still requires manual effort and expertise to understand the domain as well as the schema to write the mapping by hand.

In this thesis, we investigate to what extent an LLM agent can automate this step by generating SPARQL CONSTRUCT queries that transform a generic KG, obtained by a direct mapping, into a semantic KG\@. We extend the agent framework GRASP with a generation agent that explores the generic KG through a SPARQL endpoint, designs the model, and tests each query on the data, and an independent review agent that checks the submitted queries against the data, supported by a mechanically computed source coverage. 
% The queries depend only on the schema, so they can be inspected, corrected, and executed on the complete data.

We compare our approach with a baseline approach using a general-purpose coding agent, OpenCode, which receives the same generic KG as a file and uses the same LLM\@. We evaluate both on eight datasets: IMDb, Mondial, and MusicBrainz, with 7 to 239 tables, and five databases of the BIRD-SQL benchmark, which we model with and without competency questions. The evaluation covers on three levels: whether the output CONSTRUCT queries materialize a valid KG from the generic KG, an LLM judge that scores the model against nine criteria, and a manual analysis based on how well each KG answers a set of analytical questions, each written as a SPARQL query.

GRASP consistently produces CONSTRUCT queries that materialize a valid KG without any manual intervention, while half of the query sets of the baseline approach required fixes before they could be executed. On every dataset, the LLM judge scores the best output of GRASP higher than the best output of the baseline. Competency questions mainly help the baseline, while GRASP reaches its best results on average without them. On IMDb, its graph is also simpler and faster to query for every analytical question. GRASP scores the lowest on MusicBrainz, the dataset with the largest schema, while the baseline could not produce a result for the same dataset. These results show that an LLM agent that explores and tests against the data can produce reliable CONSTRUCT queries for a well-modeled KG, which an expert can inspect and refine for the intended use, although the quality of the model decreases for very large schemas.
